% =====================================================================
%  Abstract
% =====================================================================
\begin{abstract}
\noindent
La distribuzione quantistica di chiavi (QKD) permette a due parti di
concordare una chiave segreta con una sicurezza fondata sulle leggi della
fisica anziché sulla difficoltà computazionale. Portarla dal laboratorio
alla rete di un operatore richiede però di gestirla come qualunque altra
risorsa di rete: scoprirne la topologia, monitorarne lo stato, allocarne la
capacità ai servizi e coordinarla con le reti classiche. Il gruppo ETSI ISG
QKD ha affrontato il problema per gradi. La GS~QKD~002 individua i casi
d'uso e i requisiti di integrazione; la GS~QKD~015 definisce il modello con
cui un controller SDN vede e configura i nodi QKD di un dominio; la
GS~QKD~018 definisce l'interfaccia tra un \emph{SDN orchestrator} e il
controller di un dominio QKD, abilitando il provisioning end-to-end tra più
domini e tra rete QKD e rete di trasporto ottico. Questo lavoro ripercorre
la catena in modo ragionato: descrive le componenti di una rete SD-QKD e la
struttura di un modello YANG, analizza in dettaglio la GS~QKD~018 (casi
d'uso, modello informativo, moduli di topologia, inventario e connettività,
workflow multi-dominio, sicurezza), ne mostra l'analogia con il framework
ACTN dell'IETF (RFC~8453) e la confronta con tre esperienze reali: il
testbed multi-vendor KOREN di SK Telecom, un prototipo gerarchico basato su
ONOS sviluppato al Politecnico di Milano e la rete MadQCI di Madrid. Si
conclude con i limiti attuali dello standard e le direzioni di ricerca
aperte.

\medskip
\noindent\textbf{Parole chiave:} QKD, SDN, orchestrazione, ETSI GS QKD 015,
ETSI GS QKD 018, YANG, RESTCONF, multi-dominio, gestione delle chiavi.
\end{abstract}
